\title{Group Sequence Policy Optimization}

\begin{abstract}
We present Group Sequence Policy Optimization (GSPO), a reinforcement learning framework designed for training large language models with enhanced stability, efficiency, and effectiveness. Departing from conventional methods that utilize token-level importance ratios, GSPO formulates the importance ratio with respect to entire sequences, thereby implementing reward calculation, optimization, and clipping at the sequence level. Our results demonstrate that GSPO outperforms the GRPO algorithm both in training efficiency and overall outcomes, provides consistent stability when applied to Mixture-of-Experts (MoE) RL training, and offers pathways to streamline RL infrastructure. The advantages provided by GSPO have played a significant role in driving the advances observed in recent Qwen3 models.
\end{abstract}

\section{Introduction}
\label{sec:intro}

Reinforcement learning (RL) has become a fundamental approach for advancing the scalability of language models \citep{o1,dpsk-r1,qwq32b,qwen3}. Utilizing extensive RL techniques enables these models to acquire the proficiency necessary for addressing complex challenges, including those found in competition mathematics and programming, as it facilitates more intricate and prolonged reasoning.

A primary condition for effectively scaling RL by increasing computational resources is ensuring that training remains both stable and reliable. Nonetheless, leading RL approaches such as GRPO \citep{grpo} face pronounced challenges regarding stability when applied to the training of large-scale language models, frequently leading to catastrophic failures and irrecoverable model collapse \citep{qwen3, minimax-m1}. Such instability poses a significant barrier to further enhancing language model performance via sustained RL optimization.

In this study, we demonstrate that GRPO's instability arises primarily due to a core flaw in its handling and implementation of importance sampling weights. The improper use of these weights leads to substantial variance in training, which becomes more pronounced as the response sequence extends. Additionally, the clipping process serves to exacerbate this issue, resulting in the eventual breakdown of the model.

In response to these fundamental constraints, we introduce \textbf{Group Sequence Policy Optimization (GSPO)}, an RL framework tailored for training large language models. GSPO's primary contribution is its theoretically principled approach to defining the importance ratio using sequence likelihood \citep{click}, consistent with the foundation of importance sampling. Furthermore, GSPO derives normalized rewards by evaluating the advantages across several responses per query, thereby maintaining coherence between sequence-level reward assignment and the optimization objective.

Our experimental results reveal that GSPO markedly outperforms GRPO in terms of stability, sample efficiency, and overall training effectiveness. Notably, GSPO naturally overcomes the common instability issues found in RL training for large Mixture-of-Experts (MoE) architectures, thereby removing the necessity for elaborate stabilization techniques and streamlining the RL framework. The strengths of GSPO have played a pivotal role in realizing the substantial performance gains observed in the latest Qwen3 models. We anticipate that GSPO will serve as a resilient and scalable basis for further progress in large-scale RL training for language models.

\section{Preliminaries}

\paragraph{Notation}
Throughout this work, we represent an autoregressive language model with parameters $\theta$ as a policy $\pi_\theta$. The symbol $x$ refers to a query, while $\mathcal{D}$ stands for the set of queries. When the model generates a response $y$ for a query $x$, its probability under the policy $\pi_\theta$ is given by $\pi_\theta(y | x) = \prod_{t=1}^{|y|} \pi_\theta(y_t | x, y_{<t} )$, with $|y|$ being the length of $y$ in tokens. Each pair $(x, y)$—composed of a query and its corresponding response—can be assigned a score by a verifier $r$, producing a reward value $r(x, y)$ that falls within the interval $[ 0, 1 ]$.

\paragraph{Proximal Policy Optimization (PPO)} Proximal Policy Optimization (PPO) \citep{ppo} leverages data collected from the previous policy $\pi_{\theta_\text{old}}$ and restricts policy updates to stay close to the existing policy by utilizing a clipping approach. The method optimizes policies according to the objective below, deliberately excluding the KL regularization term for clarity, as it does not pertain to our main discussion:

\begin{align}
\mathcal{J}_\text{PPO}(\theta) = \mathbb{E}_{ x \sim \mathcal{D},\, y \sim \pi_{\theta_\text{old}}( \cdot | x) }
\left[ \frac{1}{|y|} \sum_{t=1}^{|y|} 
\min \left( w_{t}(\theta) \widehat{A}_{t},  \, \mathrm{clip} \left( w_{t}(\theta), 1 - {\varepsilon}, 1 + {\varepsilon}\right) \widehat{A}_{t} \right)
\right],
\end{align}

Here, the importance ratio for the token $y_t$ is given by $
w_{t}(\theta) = \frac{ \pi_{\theta} (y_{t} | x, y_{<t}) }{ \pi_{\theta_\text{old}} (y_{t} | x, y_{<t})}
$, where $\widehat{A}_{t}$, representing the advantage of $y_t$, is computed using an alternative value estimator, and $\varepsilon$ denotes the range used for clipping the importance ratios.

A principal difficulty with PPO in real-world applications is its significant dependence on the value model. In particular, because the value model is typically as large as the policy model, it substantially increases memory requirements and computational overhead. Additionally, PPO’s performance is tightly coupled with the accuracy of the value function estimation. Although developing an accurate value model is already a nontrivial problem, achieving scalability for longer output sequences and more sophisticated tasks amplifies these challenges further.

\paragraph{Group Relative Policy Optimization (GRPO)} As proposed by GRPO \citep{grpo}, this approach eliminates the dependence on the value model by assessing the relative advantage among a set of responses generated for an identical prompt. In particular, GRPO seeks to maximize the following objective:

\begin{align}
\label{equ:grpo}
\mathcal{J}_\text{GRPO}(\theta) = \mathbb{E}_{ x \sim \mathcal{D},\, \{y_i\}_{i=1}^G \sim \pi_{\theta_\text{old}}( \cdot | x) }
\left[ \frac{1}{G} \sum_{i=1}^{G} \frac{1}{|y_i|} \sum_{t=1}^{|y_i|} 
\min \left( w_{i,t}(\theta) \widehat{A}_{i,t},  \, \mathrm{clip} \left( w_{i,t}(\theta), 1 - {\varepsilon}, 1 + {\varepsilon}\right) \widehat{A}_{i,t} \right)
\right],
\end{align}

where $G$ is the number of generated responses to each query $x$ (i.e., the group size), and the importance ratio $w_{i,t}(\theta)$ and advantage $\widehat{A}_{i,t}$ of token $y_{i,t}$ are:

\begin{align}
    w_{i,t}(\theta)=\frac{ \pi_{\theta} (y_{i,t} | x, y_{i,<t}) }{ \pi_{\theta_\text{old}} (y_{i,t} | x, y_{i,<t})},\quad\ 
    \widehat{A}_{i,t} = \widehat{A}_{i} = \frac{r(x, y_i) - \mathrm{mean} \left( \{ r(x, y_i) \}_{i=1}^G \right) }{ \mathrm{std} \left( \{ r(x, y_i) \}_{i=1}^G \right) },
\end{align}

respectively, where all the tokens in $y_i$ share the same advantage as $\widehat{A}_{i}$.

\section{Motivation}
\label{sec:motivation}

As model sizes increase, incorporating greater sparsity—such as that seen in Mixture-of-Experts architectures—and generating longer responses, it becomes necessary to adopt large rollout batch sizes to achieve optimal hardware throughput during reinforcement learning. For the sake of sample efficiency, these sizable rollout batches are routinely divided into several smaller mini-batches for the purpose of gradient updates. However, this process inherently creates an off-policy training scenario: the generated responses $y$ are drawn according to a previous policy $\pi_{\theta_\text{old}}$, as opposed to the current, target policy $\pi_\theta$. This circumstance underpins the use of clipping strategies in algorithms like PPO and GRPO, which safeguard against incorporating samples that have drifted too far from the target policy in gradient estimation.

Although techniques such as clipping seek to control the off-policy mismatch, we observe a deeper limitation in GRPO: \textit{the objective itself is fundamentally ill-posed}. This issue is especially pronounced when large-scale models are trained on tasks requiring extended responses, often resulting in severe model degeneration. The root cause of this ill-posedness lies in an incorrect utilization of importance sampling weights within the GRPO objective. Typically, importance sampling is designed to compute the expectation of a function $f$ with respect to a target distribution $\pi_\text{tar}$, while making use of samples generated from a behavior distribution $\pi_\text{beh}$ by applying suitable re-weighting:

\begin{align}
\mathbb{E}_{z \sim \pi_\text{tar}} \left[ f(z) \right] = \mathbb{E}_{z \sim \pi_\text{beh}} \left[ \frac{ \pi_\text{tar} (z) }{ \pi_\text{beh} (z)} \, f(z) \right].
\end{align}

Importantly, the effectiveness of the importance weighting term $\frac{ \pi_\text{tar} (z) }{ \pi_\text{beh} (z)}$ in addressing distributional discrepancies depends on taking an average over many samples ($N \gg 1$) drawn from the behavior distribution $\pi_\text{beh}$.

Conversely, GRPO incorporates the importance weight $\frac{ \pi_{\theta} (y_{i,t} | x, y_{i,<t}) }{ \pi_{\theta_\text{old}} (y_{i,t} | x,y_{i,<t})}$ at each token position $t$. This weighting scheme, which relies on a single sampled token $y_{i,t}$ from the successive next-token distribution $\pi_{\theta_\text{old}}( \cdot | x, y_{i, <t})$, does not adequately address the challenge of correcting the distributional shift. Rather, it tends to inject substantial variance into the gradient estimates, with this noise compounding throughout longer generated sequences, and further intensified by the use of clipping. Our experiments indicate that such a training dynamic can precipitate a catastrophic collapse of the model—a state from which recovery proves exceptionally difficult. Following such failure, attempts to restart or restore progress, whether by resetting to earlier checkpoints, finely adjusting hyperparameters such as clipping thresholds, extending the output length, or altering RL prompts, consistently fail to revive effective training.

The observation detailed above indicates a central flaw in the architecture of GRPO. Specifically, the ineffectiveness of the token-level importance weights highlights an essential principle: \textbf{the granularity of the optimization objective must be aligned with the granularity of the reward signal}. As rewards are assigned to whole sequences, implementing off-policy corrections at the token level introduces challenges. Consequently, we are driven to abandon the token-based approach, opting instead to leverage importance weights and optimize at the \textit{sequence level} directly.

\section{Algorithm}

\subsection{GSPO: Group Sequence Policy Optimization}

Although the token-level importance weight $\frac{ \pi_{\theta} (y_{i,t} | x, y_{i,<t}) }{ \pi_{\theta_\text{old}} (y_{i,t} | x,y_{i,<t})}$ presents challenges within GRPO, we note that, specifically for language generation tasks, the \textit{sequence-level} importance weight $\frac{ \pi_{\theta} (y | x) }{ \pi_{\theta_\text{old}} (y | x)}$ carries an unambiguous theoretical interpretation. This measure quantifies the extent to which a response $y$, drawn from $\pi_{\theta_\text{old}} (\cdot | x)$, diverges from $\pi_{\theta} (\cdot | x)$. As such, it coheres with the reward structure defined at the sequence level, while simultaneously offering a principled foundation for the operation of the clipping mechanism.

Based on this straightforward observation, we propose the \textbf{Group Sequence Policy Optimization (GSPO)} algorithm. GSPO employs the following sequence-level optimization objective:

\begin{align}
\mathcal{J}_\text{GSPO} (\theta) =
\mathbb{E}_{ x \sim \mathcal{D},\, \{y_i\}_{i=1}^G \sim \pi_{\theta_\text{old}}( \cdot | x) }
\left[ 
\frac{1}{G} \sum_{i=1}^{G}
\min \left( s_{i}(\theta)  \widehat{A}_{i},\ \mathrm{clip} \left( s_{i}(\theta), 1 - {\varepsilon}, 1 + {\varepsilon} \right) \widehat{A}_{i} \right) 
\right],
\label{equ:gspo}
\end{align}

The objective function $\mathcal{J}_\text{GSPO} (\theta)$ is defined as the expectation over inputs $x$ drawn from distribution $\mathcal{D}$ and sets of samples $\{y_i\}_{i=1}^G$ generated using the old policy $\pi_{\theta_\text{old}}(\cdot|x)$. For each group of samples, the formulation calculates the average across $G$ elements, where each term involves the minimum between $s_{i}(\theta) \widehat{A}_{i}$ and $\mathrm{clip}(s_{i}(\theta), 1-{\varepsilon}, 1+{\varepsilon}) \widehat{A}_{i}$.

where we adopt the group-based advantage estimation:

\begin{align}
\widehat{A}_{i} = \frac{r(x, y_i) - \mathrm{mean} \left( \{ r(x, y_i) \}_{i=1}^G \right) }{ \mathrm{std} \left( \{ r(x, y_i) \}_{i=1}^G \right) },
\end{align}

and define the importance ratio $s_{i}(\theta)$ based on sequence likelihood \citep{click}:

\begin{align}
s_{i}(\theta) = \left( \frac{ \pi_{\theta} (y_i | x) }{ \pi_{\theta_\text{old}} (y_i | x)} \right)^{\frac{1}{|y_i|}}
=
\exp \left( \frac{1}{|y_i|} \sum_{t=1}^{|y_i|} \log \frac{ \pi_{\theta} (y_{i,t} | x, y_{i,<t}) }{ \pi_{\theta_\text{old}} (y_{i,t} | x,y_{i,<t})} \right).
\end{align}

Consequently, GSPO utilizes response-level clipping rather than token-level clipping, aiming to filter out excessively ``off-policy'' examples during the computation of gradients. This approach is consistent with the sequence-level nature of both reward assignment and optimization. Additionally, we employ length normalization in $s_{i}(\theta)$ to mitigate variance and to standardize $s_{i}(\theta)$ within a consistent numeric interval. Without this normalization, shifts in the probability of just a few tokens could cause considerable volatility in the sequence-level importance ratios, necessitating divergent clipping thresholds for responses of different lengths. It is important to highlight that, owing to distinct definitions for importance ratios, the magnitude of the clipping range used in GSPO generally differs substantially from those in prior algorithms (such as GRPO).

\subsection{Gradient Analysis}
\label{subsec:gradient}

We can derive the gradient of the GSPO objective as follows (clipping is omitted for brevity):

\begin{align}
\nabla_{\theta} \mathcal{J}_\text{GSPO} (\theta)
=&\ 
\nabla_{\theta} \mathbb{E}_{ x \sim \mathcal{D},\, \{y_i\}_{i=1}^G \sim \pi_{\theta_\text{old}}( \cdot | x) }
\left[ 
\frac{1}{G} \sum_{i=1}^{G} s_{i}(\theta) \widehat{A}_{i}
\right] \\
= &\ 
\mathbb{E}_{ x \sim \mathcal{D},\, \{y_i\}_{i=1}^G \sim \pi_{\theta_\text{old}}( \cdot | x) }
\left[ 
\frac{1}{G} \sum_{i=1}^{G}
s_{i}(\theta) \widehat{A}_{i}
\cdot \nabla_{\theta} \log s_{i}(\theta)
\right] \\
=&\ 
\mathbb{E}_{ x \sim \mathcal{D},\, \{y_i\}_{i=1}^G \sim \pi_{\theta_\text{old}}( \cdot | x) }
\left[ 
\frac{1}{G} \sum_{i=1}^{G} 
\left( \frac{ \pi_{\theta} (y_i | x) }{ \pi_{\theta_\text{old}} (y_i | x)} \right)^{\frac{1}{|y_i|}} \widehat{A}_{i} 
\cdot \frac{1}{|y_i|} \sum_{t=1}^{|y_i|} \nabla_{\theta} \log \pi_{\theta} (y_{i,t} | x, y_{i,<t}) 
\right].
\label{equ:gspo_grad}
\end{align}

For comparison, the gradient of the GRPO objective is as follows (note that $\widehat{A}_{i,t} = \widehat{A}_{i}$):

\begin{align}
\nabla_{\theta} \mathcal{J}_\text{GRPO} (\theta) 
=&\
\nabla_{\theta} \, \mathbb{E}_{ x \sim \mathcal{D},\, \{y_i\}_{i=1}^G \sim \pi_{\theta_\text{old}}( \cdot | x) }
\left[
\frac{1}{G} \sum_{i=1}^{G} \frac{1}{|y_i|} \sum_{t=1}^{|y_i|} 
w_{i,t}(\theta) \widehat{A}_{i,t}
\right] \\
=&\
\mathbb{E}_{ x \sim \mathcal{D},\, \{y_i\}_{i=1}^G \sim \pi_{\theta_\text{old}}( \cdot | x) }
\left[
\frac{1}{G} \sum_{i=1}^G \widehat{A}_i
\cdot \frac{1}{|y_i|} \sum_{t=1}^{|y_i|} 
\frac{ \pi_{\theta} (y_{i,t} | x, y_{i,<t}) }{ \pi_{\theta_\text{old}} (y_{i,t} | x, y_{i,<t}) }
\nabla_{\theta} \log \pi_{\theta}(y_{i,t} | x, y_{i,<t})
\right].
\end{align}

The key difference between GSPO and GRPO concerns \textit{the method by which each approach weights the gradients corresponding to the log likelihoods of tokens}. GRPO assigns a weight to each token based on its designated ``importance weight'' $\frac{ \pi_{\theta} (y_{i,t} | x, y_{i,<t}) }{ \pi_{\theta_\text{old}} (y_{i,t} | x,y_{i,<t})}$. These non-uniform weights can span ranges such as $(0, 1+\varepsilon]$ (when $\widehat{A}_{i} >0$) or $[1-\varepsilon, +\infty)$ (for $\widehat{A}_{i} < 0$), and are often significant. Over the course of training, these varying weights can accumulate and potentially introduce instability or unpredictable effects. By contrast, GSPO treats all tokens in a response identically with equal weights, thereby removing this source of instability present in GRPO.

\subsection{GSPO-token: A Token-level Objective Variant}

In certain settings such as multi-turn RL, a more detailed adjustment of advantages than at the overall sequence level may be required. To address this, we propose a token-level variant of the GSPO objective, referred to as \textbf{GSPO-token}, which enables advantage modification at the granularity of individual tokens:

\begin{align}
\mathcal{J}_\text{GSPO-token}(\theta)
=
\mathbb{E}_{ x \sim \mathcal{D},\, \{y_i\}_{i=1}^G \sim \pi_{\theta_\text{old}}( \cdot | x) }
\Bigg\{
\frac{1}{G} \sum_{i=1}^{G} \frac{1}{|y_i|} \sum_{t=1}^{|y_i|} 
\min \Big( s_{i,t}(\theta) \widehat{A}_{i,t},\ 
\mathrm{clip}\big(s_{i,t}(\theta),\ 1 - {\varepsilon},\ 1 + {\varepsilon}\big) \widehat{A}_{i,t} \Big)
\Bigg\}
,
\label{equ:gspo-token}
\end{align}

where

\begin{align}
s_{i,t}(\theta) 
= \mathrm{sg} \left[ s_{i}(\theta) \right]  \cdot \frac{ \pi_{\theta} (y_{i,t} | x, y_{i,<t}) }{ \mathrm{sg} \left[ \pi_{\theta} (y_{i,t} | x, y_{i,<t}) \right] },
\end{align}

and $\mathrm{sg}[\cdot]$ denotes only taking the numerical value but stopping the gradient, corresponding to the \texttt{detach} operation in PyTorch. The gradient of GSPO-token can be derived as:

\begin{align}
\nabla_{\theta} \mathcal{J}_\text{GSPO-token}(\theta)
=&\ 
\nabla_{\theta} \mathbb{E}_{ x \sim \mathcal{D},\, \{y_i\}_{i=1}^G \sim \pi_{\theta_\text{old}}( \cdot | x) }
\left[ 
\frac{1}{G} \sum_{i=1}^{G}
\frac{1}{|y_i|} \sum_{t=1}^{|y_i|} 
s_{i,t}(\theta) \widehat{A}_{i,t}
\right] \\
=&\ 
\mathbb{E}_{ x \sim \mathcal{D},\, \{y_i\}_{i=1}^G \sim \pi_{\theta_\text{old}}( \cdot | x) }
\left[
\frac{1}{G} \sum_{i=1}^{G} s_i (\theta) \cdot \frac{1}{|y_i|} \sum_{t=1}^{|y_i|} 
\widehat{A}_{i,t} \frac{ \nabla_{\theta} \pi_{\theta} (y_{i,t} | x, y_{i,<t}) }{ \pi_{\theta} (y_{i,t} | x, y_{i,<t}) }
\right] \\
=&\ 
\mathbb{E}_{ x \sim \mathcal{D},\, \{y_i\}_{i=1}^G \sim \pi_{\theta_\text{old}}( \cdot | x) }
\left[
\frac{1}{G} \sum_{i=1}^{G}  \left( \frac{ \pi_{\theta} (y_i | x) }{ \pi_{\theta_\text{old}} (y_i | x)} \right)^{\frac{1}{|y_i|}}
\cdot \frac{1}{|y_i|} \sum_{t=1}^{|y_i|}
\widehat{A}_{i,t} \nabla_{\theta} \log \pi_{\theta} (y_{i,t} | x, y_{i,<t})  
\right].
\label{equ:gspo-token_grad}
\end{align}

It is important to observe that the expression $\frac{ \pi_{\theta} (y_{i,t} | x, y_{i,<t}) }{ \mathrm{sg} \left[ \pi_{\theta} (y_{i,t} | x, y_{i,<t}) \right] }$ evaluates to 1, implying that, in terms of numerical value, $s_{i,t}(\theta)$ is equivalent to $s_{i}(\theta)$. By examining Equation~\eqref{equ:gspo} alongside~\eqref{equ:gspo-token}, as well as Equation~\eqref{equ:gspo_grad} with~\eqref{equ:gspo-token_grad}, we see that both GSPO-token and GSPO coincide numerically in the optimization target, the clipping mechanism, and the theoretical gradient under the condition that the advantages assigned to all tokens within the response $y_i$ are uniform (specifically, $\widehat{A}_{i,t} = \widehat{A}_{i}$). Notably, GSPO-token, however, permits greater adaptability by allowing the advantage values to be specified individually for each token.

\section{Experiments and Discussion}

\subsection{Empirical Results}

We conduct experiments using a cold-start model that has been fine-tuned from Qwen3-30B-A3B-Base. The results are presented through training reward trajectories and model performance evaluations on several benchmarks: AIME'24 (average Pass@1 across 32 samplings), LiveCodeBench (202410-202502, average Pass@1 across 8 samplings), and CodeForces (Elo Rating). Throughout the reinforcement learning phase, each rollout batch is divided into four mini-batches for performing gradient updates. For GSPO, we assign the left and right clipping bounds in Equation~\eqref{equ:gspo} to 3e-4 and 4e-4, respectively. To establish a baseline, we include GRPO, setting its corresponding left and right clipping intervals in Equation~\eqref{equ:grpo} to 0.2 and 0.27, with these values chosen following extensive tuning for an equitable evaluation. It should be highlighted that GRPO requires utilization of the Routing Replay training approach in order to guarantee stable MoE RL convergence—a topic elaborated on in \S~\ref{subsec:routing_replay}. In contrast, \textit{GSPO eliminates the dependency on this additional training strategy}.

As depicted in Figure~\ref{fig:results}, training with GSPO remains stable over the entire course. Notably, \textbf{GSPO consistently yields enhanced performance by allocating more computational resources, systematically refreshing the query set, and prolonging generation lengths}. In addition, GSPO outperforms GRPO in training efficiency, reaching higher levels of training accuracy and benchmark outcomes while utilizing equivalent computational budgets and query volumes. Lastly, implementing GSPO in RL training for the most recent Qwen3 models further affirms its capacity to effectively unlock the scaling potential of reinforcement learning in large language models.

\subsection{Curious Observation on Clipping Fractions}

One notable difference between GSPO and GRPO lies in the level at which clipping is applied: GSPO performs clipping across entire responses, whereas GRPO targets individual tokens. As illustrated in Figure~\ref{fig:clipping}, this results in GSPO exhibiting a fraction of clipped tokens that is approximately two orders of magnitude higher than that observed in GRPO, regardless of the chosen clipping intervals. Surprisingly, even though GSPO clips many more tokens—and, as a result, uses a smaller portion of data for model updating—its training efficiency surpasses that of GRPO. This unexpected outcome, where extensive token clipping actually enhances training efficiency, suggests that GRPO's method of estimating gradients at the token level introduces substantial noise and diminishes effectiveness in utilizing training samples. Conversely, by operating at the sequence level, GSPO is able to yield more robust and meaningful gradient signals during learning.

\subsection{Benefit of GSPO for MoE Training}
\label{subsec:routing_replay}

\paragraph{Background}
In contrast to RL training of dense architectures, MoE models present specific stability issues owing to their sparse activation mechanism. Notably, when employing the GRPO algorithm, we observed that \textit{expert-activation volatility} poses a significant barrier to stable convergence in RL training for MoE models. Specifically, following one or more gradient updates, the selection of experts associated with an identical response can shift markedly. For instance, considering the 48-layer Qwen3-30B-A3B-Base model, analysis reveals that after each RL gradient step, for the same rollout instance, about 10\% of experts triggered by the new policy $\pi_{\theta}$ differ from those activated by the previous policy $\pi_{\theta_\text{old}}$.
This effect, which grows increasingly pronounced with deeper MoE architectures, results in large and erratic swings in the token-level importance ratios $w_{i,t}(\theta) = \frac{ \pi_{\theta} (y_{i,t} | x, y_{i,<t}) }{ \pi_{\theta_\text{old}} (y_{i,t} | x,y_{i,<t})}$, undermining their reliability as elaborated in \S~\ref{sec:motivation} and~\ref{subsec:gradient}, thereby impairing the expected convergence of RL training.

\paragraph{Our Previous Approach} In our earlier work, we addressed this issue by utilizing the \textbf{Routing Replay} method during training. This strategy involves storing the set of experts activated under $\pi_{\theta_\text{old}}$, and subsequently enforcing the same routing decisions in $\pi_{\theta}$ whenever we evaluate importance ratios, where $w_{i,t}(\theta) = \frac{ \pi_{\theta} (y_{i,t} | x, y_{i,<t}) }{ \pi_{\theta_\text{old}} (y_{i,t} | x,y_{i,<t})}$. By matching the activated experts between $\pi_{\theta}$ and $\pi_{\theta_\text{old}}$ for every token $y_{i,t}$, the computations for $\pi_{\theta} (y_{i,t} | x, y_{i,<t})$ and $\pi_{\theta_\text{old}} (y_{i,t} | x,y_{i,<t})$ are performed over identical subnetworks. This alignment stabilizes the importance weights on a per-token basis and supports consistent gradient-based optimization by preventing divergence in the activated computational paths. As illustrated in Figure~\ref{fig:routing_replay}, Routing Replay is a key component for achieving reliable convergence in GRPO training within Mixture of Experts models.

\paragraph{Benefit of GSPO} While Routing Replay facilitates the convergence of GRPO training in MoE models, its reliance on reapplying routing configurations comes at the cost of increased memory usage and communication overhead, potentially restricting the effective capacity of the MoE framework. Conversely, as depicted in Figure~\ref{fig:results}, GSPO does not require Routing Replay, allowing for the standard computation of importance ratios $s_i(\theta)$, thereby ensuring stable optimization and typical convergence. The primary observation is that GSPO concentrates on overall sequence likelihood (i.e., $\pi_{\theta} (y_i | x)$), demonstrating robustness with respect to fluctuations in token-level likelihoods (i.e., $\pi_{\theta} (y_{i,t} | x, y_{i,<t})$). Because the language modeling ability of the MoE model is consistently retained, abrupt changes in the sequence likelihood are unlikely. In essence, GSPO effectively mitigates the volatility associated with expert activations in MoE models, eliminating the necessity for intricate methods such as Routing Replay. As a result, the training pipeline becomes both simpler and more robust, and the model can fully exploit its available capacity without imposed limitations.

\subsection{Benefit of GSPO for RL Infrastructure}

Due to precision inconsistencies between training platforms such as Megatron and inference platforms like SGLang and vLLM, it is common practice to utilize the training engine to re-evaluate the likelihoods of generated responses according to the previous policy $\pi_{\theta_\text{old}}$. Nevertheless, since GSPO relies on sequence-level rather than token-level likelihoods during optimization, it is inherently more robust to such precision variances. As a result, GSPO enables the direct use of likelihood scores from the inference engine for optimization purposes, circumventing the necessity for additional calculations by the training engine. This approach offers significant advantages, particularly in contexts such as partial rollout, multi-turn RL, and architectures that separate training and inference.

\section{Conclusion}

We introduce Group Sequence Policy Optimization (GSPO), a novel reinforcement learning algorithm specifically designed for the training of large language models. Built upon the core ideas of importance sampling, GSPO establishes importance ratios derived from sequence likelihoods, and applies sequence-level clipping, rewarding, and optimization. When compared to GRPO, GSPO achieves markedly enhanced training stability, greater efficiency, and superior performance. Its strengths are especially evident in large-scale reinforcement learning applied to MoE models, providing the basis for the outstanding advancements realized in recent Qwen3 models. With GSPO serving as a robust and scalable algorithmic foundation, we intend to further extend RL methods and anticipate significant progress in the development of intelligence as a result.

\end{document}